\keywordscn{肌肉驱动\quad 生物力学\quad MuJoCo\quad 人机交互\quad 方向盘控制\quad 精度测试}
\begin{abstractzh}

本课题以肌肉驱动生物力学人机交互为核心，面向视觉引导下的方向盘精准操控需求，设计并实现一套基于 MuJoCo 仿真平台的交互系统。首先梳理生物力学建模、肌肉驱动控制、视觉引导与人机交互精度评估的国内外研究现状，明确现有技术在通用化建模、跨个体适配、低成本仿真等方面的不足。其次以上肢肌肉--骨骼动力学为理论基础，采用 MuJoCo 物理引擎构建人体双臂模型与方向盘物理模型，完成肌肉激活与方向盘转角的耦合驱动。然后设计自适应非线性映射算法与视觉引导实时反馈模块，实现肌肉发力到方向盘转动的平滑控制。最后通过功能测试、精度测试、稳定性测试与易用性测试，以 MAE、RMSE、响应延迟、连续运行稳定性为指标完成系统验证。实验结果表明，本系统方向盘角度跟踪平均绝对误差小于 3$^\circ$，响应延迟低于 10\,ms，连续运行 1 小时无明显漂移，能够稳定完成视觉引导下的方向盘操控任务。本研究可为驾驶模拟、康复训练、工业协同操控等场景提供轻量化、高生物适配性的人机交互方案。

\end{abstractzh}